\documentclass[10pt]{article}
\usepackage{amsmath}
\usepackage{setspace}
\onehalfspacing
\usepackage{amssymb}
\usepackage{color}
\usepackage{fullpage}
\usepackage{graphicx}
\usepackage{cancel}

\usepackage{fullpage}

\newcommand{\rhintbox}[1]{\begin{center} \singlespacing
  \framebox[.9\textwidth][c]{\begin{minipage}{.8\textwidth}\vspace{.1em}#1
\end{minipage}}\end{center}\onehalfspacing}


\newcommand{\R}{R}

\newcommand{\answer}[1]{\begin{quote} {\color{red} #1 }\end{quote}}
\newcommand{\var}{\textrm{var}}
\newcommand{\cov}{\textrm{cov}}

\title{QTM 220 Midterm Solution: Fall 2023 (Sept. 21)}
\date{}


\begin{document}

\maketitle


\subsection*{Problem 1}

Suppose we store univariate sample data in a vector $y$ of length $n$ and we run the following code:

\begin{verbatim}
B=10000
y.bar = rep(NA,B)
for(b in 1:B){y.bar[b]= mean(sample(y,size=n,replace=T))}
\end{verbatim}


\begin{itemize}
\item[a)]  What does the following code produce?

\begin{verbatim}
mean(y)
\end{verbatim}

\color{red} sample mean \color{black}

\item[b)] What does the following code produce?
\begin{verbatim}
sd(y)
\end{verbatim}

\color{red} sample standard deviation \color{black}


\item[c)]  What does the following code produce?

\begin{verbatim}
sd(y.bar)
\end{verbatim}

\color{red} Estimated standard error \color{black}

\item[d)] What does the following code produce?
\begin{verbatim}
quantile(y.bar,probs=c(.05,.95)) 
\end{verbatim}

\color{red} 90\% CI \color{black}

\item[e)] What does the following code produce?
\begin{verbatim}
mean(y) - 1.96* sd(y)/sqrt(n)
mean(y) + 1.96* sd(y)/sqrt(n)
\end{verbatim}

\color{red} 95\% CI \color{black}

\item[f)] How could you re-write the code in part e) using sd(y.bar)?

\color{red} \begin{verbatim}
mean(y) - 1.96* sd(y.bar)
mean(y) + 1.96* sd(y.bar)
\end{verbatim}
 \color{black}



\end{itemize}

\clearpage

\subsection*{Problem 2}

Suppose we store bivariate sample data in a data frame dfsam with the first column dfsam[,1] the outcome variable $y$ and the second column dfsam[,2]  the binary covariate $x$. The number of rows is $n$.

\begin{itemize}
\item[a)]  What does the following code produce?

\begin{verbatim}
mean(y[x==0]) 
\end{verbatim}

\color{red} sample mean for the $x=0$ observations \color{black}
\item[b)]  Explain why the following code will produce the same thing.

\begin{verbatim}
lm(y~x,data=dfsam)$coef[1]
\end{verbatim}

\color{red} with a binary covariate the linear regression intercept is the mean for the $x=0$ group  \color{black}



\item[c)]  What does the following code produce?

\begin{verbatim}
mean(y[x==1]) - mean(y[x==0]) 
\end{verbatim}

\color{red} difference in sample means \color{black}
\item[d)]  Explain why the following code will produce the same thing.

\begin{verbatim}
lm(y~x,data=dfsam)$coef[2]
\end{verbatim}

\color{red} with a binary covariate the linear regression slope is a difference in means \color{black}



\end{itemize}

\clearpage

%----------------------------------------------%%
% Problem 1
%----------------------------------------------%%
 \subsection*{Problem 3}

 Suppose you have a sample $Y_1, Y_2, \ldots, Y_n$ from a population of unknown distribution, with mean $\mu$ and standard deviation $\sigma$ (and therefore variance $\sigma^2$).  While you do not know the population distribution, you do know that each $Y_i$ is independent and identically distributed (i.i.d.).  \\

 Consider the following estimator for $\mu$: $\tilde{Y} = \frac{1}{n-2} \sum_{i=1}^n Y_i$.

 \begin{itemize}
 \item[a)] What is $E[\tilde{Y}]$?

 \color{red}
 \begin{align*}
 E[\tilde{Y}] &= E[ \frac{1}{n-2} \displaystyle \sum_{i=1}^n Y_i] \\
 &= \frac{1}{n-2} E[ \sum_{i=1}^n Y_i]\\
 &= \frac{1}{n-2} (\sum_{i=1}^n E[Y_i])\\
 &=\frac{1}{n-2} (\sum_{i=1}^n \mu )\\
 &=\frac{1}{n-2} n\mu \\
 &= \frac{n}{n-2}\mu
 \end{align*}
 \color{black}


\item[b)] Is $\tilde{Y}$ biased or unbiased for $\mu$?

\color{red} biased \color{black}



% \item[b)] What is the bias of $E[\tilde{Y}]$ for $\mu$?
%
% \color{red}
% \begin{align*}
% E[\tilde{Y} - \mu] &= \frac{n}{n-2}\mu - \mu \\
% &= (\frac{n}{n-2} - 1)\mu \\
% &= (\frac{n}{n-2} - \frac{n-2}{n-2})\mu \\
% &= \frac{2}{n-2}\mu
%  \\\end{align*}
% \color{black}

 \item[c)] Showing your work, calculate $Var[\tilde{Y}]$.

 \color{red}
 \begin{align*}
 V[\tilde{Y}] &= V[ \frac{1}{n-2} \displaystyle \sum_{i=1}^n Y_i] \\
 &= \frac{1}{(n-2)^2} V[ \sum_{i=1}^n Y_i]\\
 &= \frac{1}{(n-2)^2} (\sum_{i=1}^n V[Y_i])\\
 &=\frac{1}{(n-2)^2} (\sum_{i=1}^n \sigma^2 )\\
 &=\frac{1}{(n-2)^2} n\sigma^2 \\
 &= \frac{n\sigma^2}{(n-2)^2}
 \end{align*}
 \color{black}


% \item[d)] Is $\tilde{Y}$ a \emph{consistent} estimator for the population mean $\mu$? Why or why not?
% \color{red}
% Yes, both its bias and variance go to zero.
% \color{black}

 \end{itemize}


\subsection*{Problem 4}


Suppose we are studying a trait that is distributed differently in two different groups in the population. We call these groups the ``low'' group (characterized by $Y_{low}$) and the ``high'' group (characterized by $Y_{high}$). $Y_{low}$ and $Y_{high}$ have the following distributions:\\

$$Y_{low} \hspace{0.1cm} \underset{i.i.d.}{\sim} \hspace{0.1cm} N\hspace{0.1cm}(\hspace{0.1cm}\mu,\hspace{0.1cm} \sigma^2\hspace{0.1cm}) $$

$$Y_{high} \hspace{0.1cm} \underset{i.i.d.}{\sim} \hspace{0.1cm} N\hspace{0.1cm}(\hspace{0.1cm}2 \mu,\hspace{0.1cm} \sigma^2\hspace{0.1cm}) $$

\noindent The parameter of interest is $\mu$. Suppose we are able to conduct separate random samples with replacement from the groups of sizes $n_{low}$ and $n_{high}$. \\%Assume that the probability of being in the lowogenous group is 0.5 and the probability of being in the higherogeneous group is also 0.5. \\

%Consider the following 100 observations randomly sampled from a Bernoulli distribution with parameter $p$\hspace{0.05cm}:\\

%$$ Y_1, Y_2, Y_3, ..., Y_{100} \hspace{0.1cm}\underset{i.i.d.}{\sim} \hspace{0.1cm} Bern\hspace{0.1cm}(p) $$


\begin{itemize}
\item[a)]  Consider this estimator for $\mu$: $ \hat{\mu} = \sum_{j=1}^{n_{high}} a_j Y_{j,high} - \sum_{i=1}^{n_{low}} b_iY_{i,low}$ \\

Find the expected value of this estimator, $E[\hat{\mu}]$. \\

\color{red}
\begin{align*}
E[\hat{\mu}] &= E\Big[\sum_{j=1}^{n_{high}} a_j Y_{j,high} - \sum_{i=1}^{n_{low}} b_iY_{i,low}]  \\
&= E\Big[\sum_{j=1}^{n_{high}} a_j Y_{j,high}] - E\Big[\sum_{i=1}^{n_{low}} b_iY_{i,low}]  \\
&= \sum_{j=1}^{n_{high}} E\Big[a_j Y_{j,high}] - \sum_{i=1}^{n_{low}} E\Big[b_iY_{i,low}]  \\
&= \sum_{j=1}^{n_{high}} 2\mu a_j  - \sum_{i=1}^{n_{low}}\mu b_i  \\
&= \mu(2\sum_{j=1}^{n_{high}}  a_j  - \sum_{i=1}^{n_{low}} b_i ) \\
\end{align*}


\color{black}
%\item[b)] Derive a condition using $\sum_{j=1}^{n_{high}}  a_j$ and $\sum_{i=1}^{n_{low}} b_i$ that will ensure unbiasedness. \\
%\color{red}
%
%$$2\sum_{j=1}^{n_{high}}  a_j  - \sum_{i=1}^{n_{low}} b_i = 1$$
%
%\color{black}

\item[bonus)] Showing your work, calculate $Var[\hat{\mu}]$.

\color{red}
\begin{align*}
V[\hat{\mu}] &= V\Big[\sum_{j=1}^{n_{high}} a_j Y_{j,high} - \sum_{i=1}^{n_{low}} b_iY_{i,low}]  \\
&= V\Big[\sum_{j=1}^{n_{high}} a_j Y_{j,high}] + V\Big[\sum_{i=1}^{n_{low}} b_iY_{i,low}]  \\
&= \sum_{j=1}^{n_{high}} V\Big[a_j Y_{j,high}] + \sum_{i=1}^{n_{low}} V\Big[b_iY_{i,low}]  \\
&= \sigma^2 \sum_{j=1}^{n_{high}} a_j^2  + \sigma^2 \sum_{i=1}^{n_{low}} b_i^2  \\
&= \sigma^2(\sum_{j=1}^{n_{high}}  a_j^2  + \sum_{i=1}^{n_{low}} b_i^2 ) \\
\end{align*}


\color{black}


%\item[d)] Under what conditions is $\hat{\mu}$ a \emph{consistent} estimator for $\mu$?
%
%\item[e)] Provide some example values of $a_j$  and $b_i$ that satisfy your conditions in parts (b) and (d) and discuss the intuition behind the implied estimators.
%
\end{itemize}

\clearpage




\end{document}
